\documentclass{ximera}
\title{Regular Polygons}

\newcommand{\pskip}{\vskip 0.1 in}

\begin{document}
\begin{abstract}
Using vectors to construct a regular polygon.
\end{abstract}
\maketitle


\begin{question} \label{Qc56r4322}
The points $A$ and $B$, with respective coordinates $(1,2)$ and $(5,1)$, are consecutive vertices of a regular pentagon $ABCDE$, oriented counterclockwise as shown below.


\pdfOnly{Access Desmos interactive at \href{https://www.desmos.com/calculator/77xkbgwqvb}{https://www.desmos.com/calculator/77xkbgwqvb}}

 
\begin{onlineOnly}
    \begin{center}
\desmos{77xkbgwqvb}{900}{600}
\end{center}
\end{onlineOnly}

\href{https://www.desmos.com/calculator/77xkbgwqvb}{Regular Pentagon}

\begin{enumerate}
\item Use vectors to find the exact coordinates of the third vertex $C$.

\item Check your work by inputting these coordinates in Line 39 of the worksheet.

\item \emph{Challenge:} Use vectors to find the exact coordinates of the two other vertices.

\end{enumerate}

\emph{Hint:} Activate the folder in Line 2 for a hint. Start by determining the measure of the angle $\alpha$. Then find components of the perpendicular vectors ${\bf u}$ and ${\bf v}$.

\begin{explanation}
The idea is to find the components of the vector $\overrightarrow{BC}$. Then we can use these to find the components of the vector
\[
   \overrightarrow{OC} = \overrightarrow{OB} + \overrightarrow{BC}
\]
and hence the coordinate of $C$.

To find the components of the vector $\overrightarrow{BC}$ we start with the vector
\[
  {\bf u} = \overrightarrow{AB}  = \langle \answer{4}, \answer{-1}  \rangle. 
\]
Then rotate this vector $\pi/2$ radians counterclockwise to get the vector
\[
  {\bf v} = \langle \answer{1} , {4}\rangle . 
\] 
Then 
\[
  \overrightarrow{BC} = (\cos \alpha) {\bf u} + (\sin\alpha) {\bf v} ,
\]
where $\alpha$ is the counterclockwise angle from ${\bf u}$ to $\overrightarrow{BC}$.

To finish the problem from here, we need to find the measure of the angle $\alpha$. This is one of five equal exterior angles of the regular pentagon. Since in walking one around the pentagon, you turn through these five angles and end up facing the direction in which you started, these angles sum to $2\pi$ radians. So each has measure $\alpha =2\pi/5$. I'll let you finish the problem from here.

\end{explanation}

\end{question}

\end{document}